% from ega4-21a-fr.tex
\oldpage[IV-4]{269}

\begin{definition}[21.5.5]
\EGAReaderTarget{ega4.c04.def.21-5-5}{IV.21.5.5}
Étant donné un morphisme fini $f:X'\to X$ de préschémas, vérifiant l'une des conditions I), II) de \sref{IV.21.5.3}, on appelle \emph{image directe} (ou \emph{norme}) d'un diviseur $D'$ sur $X'$ par $f$, et l'on note $f_*(D')$ (ou $\mathrm{N}_{X'/X}(D')$), le diviseur sur $X$ qui correspond canoniquement \sref{IV.21.2.11} à la classe du couple $(\mathrm{N}_{X'/X}(\sh{O}_{X'}(D')),\mathrm{N}_{X'/X}(s_{D'}))$.
\end{definition}

Il résulte aussitôt de cette définition, compte tenu de \sref{IV.21.2.9.4} et de \sref{IV.21.5.3.2}, que si $D_1'$, $D_2'$, $D'$ sont des diviseurs sur $X'$, on a
\[
\label{IV.21.5.5.1}
  f_*(D_1'+D_2')=f_*(D_1')+f_*(D_2')
  \tag{21.5.5.1}
\]
et $D'\geq0$ entraîne $f_*(D')\geq0$, autrement dit, $D'\mapsto f_*(D')$ est un homomorphisme croissant du groupe ordonné $\operatorname{Div}(X')$ dans le groupe ordonné $\operatorname{Div}(X)$.
La définition \sref{IV.21.5.5} montre d'ailleurs aussitôt que pour tout ouvert $U$ de $X$, on a $f_*^U(D'|f^{-1}(U))=f_*(D')|U$ ($f^U$ étant la restriction $f^{-1}(U)\to U$ de $f$), et les homomorphismes $f_*^U$, pour $U$ variable, définissent donc un homomorphisme de faisceaux de groupes ordonnés
\[
\label{IV.21.5.5.2}
  \mathrm{N}_{X'/X}:f_*(\sh{D}iv_{X'})\to\sh{D}iv_X.
  \tag{21.5.5.2}
\]
En outre, pour tout ouvert $U$ de $X$, tout $\sh{O}_{X'}$-Module inversible $\sh{L}'$ et toute section méromorphe régulière $s'$ de $\sh{L}'$ au-dessus de $f^{-1}(U)$, on a, d'après les définitions précédentes et \sref{IV.21.1.4}
\[
\label{IV.21.5.5.3}
  \operatorname{div}_U(\mathrm{N}(s'))=f_*^U(\operatorname{div}_{f^{-1}(U)}(s')).
  \tag{21.5.5.3}
\]

\begin{proposition}[21.5.6]
\label{IV.21.5.6}
Soit $f:X'\to X$ un morphisme fini localement libre et supposons que $f_*(\sh{O}_{X'})$ soit de rang constant $n$.
Alors, pour tout diviseur $D$ sur $X$, $f_*(D)$ est défini et l'on a
\[
\label{IV.21.5.6.1}
  f_*(f^*(D))=nD.
  \tag{21.5.6.1}
\]
\end{proposition}

\begin{proof}
La première assertion résulte de ce que $f$ est plat \sref{IV.21.4.5}, et la seconde est conséquence immédiate des définitions et de (\textbf{II}, 6.5.3.2).
\end{proof}

\begin{proposition}[21.5.7]
\label{IV.21.5.7}
Soient $f:X'\to X$ un morphisme fini vérifiant l'une des hypothèses I), II) de \sref{IV.21.5.3}, $f':X''\to X'$ un morphisme fini localement libre de rang constant $n$.
Alors $f''=f\circ f':X''\to X$ vérifie la même hypothèse que $f$, et pour tout diviseur $D''$ sur $X''$, on a
\[
\label{IV.21.5.7.1}
  f_*''(D'')=f_*(f_*'(D'')).
  \tag{21.5.7.1}
\]
\end{proposition}

Compte tenu de la définition \sref{IV.21.5.5}, il suffit de prouver le résultat suivant :

\begin{lemma}[21.5.7.2]
\label{IV.21.5.7.2}
Sous les hypothèses de \sref{IV.21.5.7}, on a un isomorphisme fonctoriel
\[
\label{IV.21.5.7.3}
  \mathrm{N}_{X''/X}(\sh{L}'')\xrightarrow{\sim}\mathrm{N}_{X'/X}(\mathrm{N}_{X''/X'}(\sh{L}''))
  \tag{21.5.7.3}
\]
dans la catégorie des $\sh{O}_X$-Modules inversibles.
\end{lemma}

En effet, compte tenu de la définition de la norme d'une section de $\sh{L}''$ (\textbf{II}, 6.5.3) et de la définition \sref{IV.21.5.5}, on obtiendra aussitôt \sref{IV.21.5.7.1}.
Pour prouver \sref{IV.21.5.7.2}, il suffit, compte tenu des définitions de (\textbf{II}, 6.5.2 et 6.5.3), de prouver que pour toute section $s$ de $f_*''(\sh{O}_{X''})=f_*(f_*'(\sh{O}_{X''}))$ au-dessus d'un ouvert $U$ de $X$, on a
\[
\label{IV.21.5.7.4}
  \mathrm{N}_{f_*''(\sh{O}_{X''})|\sh{O}_X}(s)=\mathrm{N}_{f_*(\sh{O}_{X'})|\sh{O}_X}(\mathrm{N}_{f_*'(\sh{O}_{X''})|\sh{O}_{X'}}(s)).
  \tag{21.5.7.4}
\]

